% Nouveau EGA V (cote 25): archive cote 25, site page 24 (PDF page 25).
% typescript. Draft compared: none.
% AI-transcribed by Claude Opus 5.5 (claude-opus-5-5, adaptive thinking, effort high), Message Batches API, batch ega25b (msgbatch_01A7n5YuXubUPNK1zsepi8RG), 2026-11-13; not yet checked by eye.
\scanpage{25/24}
\mspage{\hand{21}}

satisfaite notamment si on supposer $F_s$ Cohen-Macaulay (à fortiori, si
F est CM sur S); de plus, dans ce cas, $G_s$ est CM (car localement il se
déduit de $F_{k(s)}$ qui l'est en divisant par un $\phi . F_{k(s)}$ , \deleted{avec} \inserted{où} $\phi$ \inserted{est} $F_{k(s)}$-%
régulier). Ces mêmes remarques pouvaient \inserted{(et devaient)} d'ailleurs se faire localement
en haut, pour caractériser les points de U (au lieu de ceux de V).

Utilisant maintenant 12.1.1. et 12.1.4. , on obtient:
\uline{Théorème} 5.6. Supposons F de présentation finie, plat rel. à S. Soit
P l'une des propriétés \xout{\unsure{satisfaites pour un module de type fini M sur}}
\xout{\unsure{un anneau local noethérien A sur un corps K}} (i) à (viii) de 12.1.1. ,
ou \xout{\unsure{(ii) de exposé M (réf)}} (dans le cas où on suppose $F=\underline{O}_X$) l'une des
\xout{\unsure{xxxx(ii) M satisfait $(S_k)$ (réf)}} propriétés (i) à (iv)\hand{/} de 12.1.4. Soit
\margin{\unsure{1r}} $U_P$ \xout{\unsure{(iii) M est Cohen-Macaulay}} l'ensemble des $x \in X_{\hand{\check{}}P}$ tels que , si \hand{$\zeta$}
\xout{xxx(iv)} désigne l'image de x dans $\hand{\check{P}}$ , la propriété P soit satisfaite
par $G_{\hand{\zeta}}$ (resp. \xout{\unsure{0}} $Y_{\hand{\zeta}}$ ) en le point x, \uline{et} que $\phi_{\hand{\zeta}}$ soit $F_{k(\hand{\zeta})}$-régulière en x.
Alors $U_P$ est ouvert, et $\underline{G}\hand{|}U_P$ est plat relativement à S.

En effet, par définition même, on a $U_P \subset U$ (notation de 5.3. a)),
et on applique par.12 à $U \to \hand{\check{P}}$ et $F_{\hand{\check{P}}}\hand{|}U$.
\uline{Corollaire} 5.7. Suppososons F de présentation finie, plat rel à S,
et supp F propre sur S (p.ex. X propre sur S). Soit $V_P$ l'ensemble des
$\hand{\zeta\in}\hand{\check{P}}$ tels que $G_{\hand{\zeta}}$ (resp. $Y_{\hand{\zeta}}$ ) satisfasse à la propriété P, et que $\phi_{\hand{\zeta}}$ soit
$F_{k(\hand{\zeta})}$-régulière. Sous ces conditions, $V_P$ est ouvert, (et il est d'ailleu
rs constructible en tous cas, i.e. sans hypothèse de platitude ni de
\unsure{prppreté|propreté}).

Il me semble que du point de vue présentation, on ne peut laisser
5.6. tel quel, avec une simple référence à des conditions énumérées dans
un autre fascicule; mais il faudra refaire une liste explicite (i) (ii)
.... des propriétés qu'on a en vue. De plus, signaler \inserted{(dans 5.1. peut-être)} que le cas P= géom.
normal (avec S =Spec(k) pour sur) est dû à \uline{Seidenberg}.

\margin{24}
